\documentclass[11pt,a4paper]{article}
\usepackage[margin=25mm,headheight=16pt]{geometry}
\usepackage{fontspec,kotex,xcolor,hyperref,fancyhdr}
\setmainfont{DejaVu Serif}
\setmainhangulfont{Noto Sans CJK KR}
\definecolor{navy}{HTML}{173E59}
\pagestyle{fancy}\fancyhf{}
\fancyhead[L]{\small GuardSynth CoC}
\fancyhead[R]{\small 제목·초록 검토안 · 한글}
\fancyfoot[L]{\small v0.1 · 2026-09-30}
\fancyfoot[R]{\small\thepage}
\renewcommand{\headrulewidth}{0.3pt}
\setlength{\parindent}{0pt}
\setlength{\emergencystretch}{3em}
\linespread{1.15}
\hypersetup{pdftitle={GuardSynth: CoC 기반 VLM 학습을 위한 행동 제약의 정형 명세와 검증},pdfsubject={Title and abstract for author review},pdfauthor={GuardSynth project}}
\begin{document}
\vspace*{8mm}
{\Large\bfseries\color{navy} GuardSynth: CoC 기반 VLM 학습을 위한 행동 제약의 정형 명세와 검증\par}
\vspace{8mm}
{\large\bfseries 초록\par}
\vspace{3mm}
Chain of Cognition(CoC)은 장면 이해와 행동 판단을 연결하여 시각언어모델(VLM)의 학습에 활용된다. 그러나 행동 제약의 적용 조건과 해제 시점이 명시되지 않은 CoC만으로 학습하면 올바른 행동 선택에 한계가 있다. 본 논문은 CoC에 명시적 행동 제약을 결합하는 GuardSynth를 제안한다. 행동 제약의 범위와 구성 요소를 정의하고, 적용 조건·금지 행동·유지 및 해제 조건·시간·근거를 Evidence-Bound Lifecycle Contracts(EBLC)로 명세한다. 이어 제한된 논리 모델에서 일관성과 지정 속성을 검증하고, 명세 필드에 대응하는 통제 자연어(Controlled Natural Language, CNL)를 생성하여 CoC 학습 입력에 추가한다. 통제된 합성 시각·시간 과제에서 Qwen3-VL-2B-Instruct를 동일한 LoRA 학습 예산과 세 개의 시드로 비교하였다. CoC 단독, 기존 자연어 제약 추가, EBLC 기반 CNL 추가 조건의 평균 행동 선택 정확도는 각각 49.86\%, 96.81\%, 98.06\%였으며, 제약 위반율은 19.17\%, 1.39\%, 0.83\%였다. 결과는 CoC 행동 제약을 체계적으로 명세·검증하고 학습 입력으로 연결하는 방법과, 명시적 제약 제공에 따른 행동 선택 개선을 보여준다. 다만 학습과 평가 모두에서 제약을 제공한 통제 과제의 결과이며, 검증 자체의 성능 개선 효과나 실제 도로 안전성을 입증하는 것은 아니다.
\end{document}
